\subsection{LC Pass Filter}

\subsubsection{Overview}
Constructed from an inductor and capacitor in series. It
is a second order filter. An LC Pass Filter can be broken
down into two types: Low Pass Filter and High Pass.
\begin{enumerate}
  \item \textbf{Low Pass Filter:} Allows low frequency
    signals to pass while attenuating high frequency signals.
  \item \textbf{High Pass Filter:} Allows high frequency
    signals to pass while attenuating low frequency signals.
\end{enumerate}

\subsubsection{Schematic}
\begin{circuitfig}
  % Paths, nodes and wires:
  \draw (4.25, 6.5) -- (5.25, 6.5);
  \node[shape=rectangle, minimum width=0.715cm, minimum height=0.465cm](N1) at (5.625, 6.5){} node[anchor=center] at (N1.text){$V_{out}$};
  \draw (2.75, 8.25) to[sinusoidal voltage source, l={$V_1$}] (2.75, 6.75);
  \draw (2.75, 4.75) to[sinusoidal voltage source, l_={$V_2$}] (2.75, 6.25);
  \node[ground, rotate=-90] at (2.75, 6.75){};
  \node[ground, rotate=-90] at (2.75, 6.25){};
  \draw (2.75, 4.75) -- (4.25, 4.75);
  \draw (2.75, 8.25) -- (4.25, 8.25);
  \node[shape=rectangle, minimum width=2.965cm, minimum height=0.465cm] at (3.5, 4.25){} node[anchor=north, align=center, text width=2.577cm, inner sep=6pt] at (3.5, 4.5){Full Schematic};
  \draw (8.625, 6.5) -- (9.625, 6.5);
  \node[shape=rectangle, minimum width=0.715cm, minimum height=0.465cm](N2) at (10, 6.5){} node[anchor=center] at (N2.text){$V_{out}$};
  \draw (7.125, 8.25) to[sinusoidal voltage source, l={$V_1$}] (7.125, 6.75);
  \node[ground, rotate=-90] at (7.125, 6.75){};
  \node[ground, rotate=-90] at (7.125, 4.75){};
  \draw (7.125, 4.75) -- (8.625, 4.75);
  \draw (7.125, 8.25) -- (8.625, 8.25);
  \node[shape=rectangle, minimum width=2.965cm, minimum height=0.465cm] at (7.875, 4.25){} node[anchor=north, align=center, text width=2.577cm, inner sep=6pt] at (7.875, 4.5){Low Pass};
  \draw (13, 6.5) -- (14, 6.5);
  \node[shape=rectangle, minimum width=0.715cm, minimum height=0.465cm](N3) at (14.375, 6.5){} node[anchor=center] at (N3.text){$V_{out}$};
  \draw (11.5, 4.75) to[sinusoidal voltage source, l_={$V_2$}] (11.5, 6.25);
  \node[ground, rotate=-90] at (11.5, 8.25){};
  \node[ground, rotate=-90] at (11.5, 6.25){};
  \draw (11.5, 4.75) -- (13, 4.75);
  \draw (11.5, 8.25) -- (13, 8.25);
  \node[shape=rectangle, minimum width=2.965cm, minimum height=0.465cm] at (12.25, 4.25){} node[anchor=north, align=center, text width=2.577cm, inner sep=6pt] at (12.25, 4.5){High Pass};
  \draw (4.25, 8.25) to[american inductor, l={$L_1$}] (4.25, 6.5);
  \draw (8.625, 8.25) to[american inductor, l={$L_1$}] (8.625, 6.5);
  \draw (13, 8.25) to[american inductor, l={$L_1$}] (13, 6.5);
  \draw (4.25, 6.5) to[capacitor, l={$C_1$}] (4.25, 4.75);
  \draw (8.625, 6.5) to[capacitor, l={$C_1$}] (8.625, 4.75);
  \draw (13, 6.5) to[capacitor, l={$C_1$}] (13, 4.75);
\end{circuitfig}

\subsubsection{Transfer Function}
Applying KCL at $V_{out}$ on Full Schematic, we get:
\[
  \frac{(V_{out} - V_{1})}{sL_1} + (V_{out} - V_{2}){sC_1} = 0
\]
\paragraph{Superposition:}
\begin{enumerate}
  \item Due to \(V_{1}\) alone (\(V_{2} = 0V\)):
    This is the low-pass characteristic of the LC Pass Filter.
    \begin{align*}
      \frac{(V_{out} - V_{1})}{sL_1} + (V_{out} - 0)sC_1 &= 0 \\
      \frac{V_{out}}{sL_1} - \frac{V_{1}}{sL_1} + V_{out}sC_1 &= 0 \\
      V_{out}(sC_1 + \frac{1}{sL_1}) - \frac{V_{1}}{sL_1} &= 0 \\
      V_{out}(sC_1 + \frac{1}{sL_1}) &= \frac{V_{1}}{sL_1} \\
      V_{out} &= \frac{\frac{V_{1}}{sL_1}}{(sC_1 + \frac{1}{sL_1})} \\
      V_{out} &= \frac{\frac{V_{1}}{sL_1}}{\frac{s^2C_1L_1 + 1}{sL_1}} \\
      V_{out} &= \frac{V_{1}}{s^2C_1L_1 + 1} \\
      H_1(s) = \frac{N_1(s)}{D_1(s)} = \frac{V_{out_1}}{V_{1}} &= \frac{1}{s^2C_1L_1 + 1}
    \end{align*}
  \item Due to \(V_{2}\) alone (\(V_{1} = 0V\)):
    This is the high-pass characteristic of the LC Pass Filter.
    \begin{align*}
      \frac{(V_{out} - 0)}{sL_1} + (V_{out} - V_{2})sC_1 &= 0 \\
      \frac{V_{out}}{sL_1} + V_{out}sC_1 - V_{2}sC_1 &= 0 \\
      V_{out}(\frac{1}{sL_1} + sC_1) &= V_{2}sC_1 \\
      V_{out} &= \frac{V_{2}sC_1}{(\frac{1}{sL_1} + sC_1)} \\
      V_{out} &= \frac{V_{2}sC_1}{\frac{s^2C_1L_1 + 1}{sL_1}} \\
      V_{out} &= \frac{V_{2}s^2C_1L_1}{s^2C_1L_1 + 1} \\
      H_2(s) = \frac{N_2(s)}{D_2(s)} = \frac{V_{out_2}}{V_{2}} &= \frac{s^2C_1L_1}{s^2C_1L_1 + 1}
    \end{align*}
\end{enumerate}
\paragraph{Total Output:}
\begin{align*}
  V_{out} &= V_{out_1} + V_{out_2} \\
  V_{out} &= H_1(s)V_{1} + H_2(s)V_{2} \\
  V_{out} &= \underbrace{\frac{1}{s^2C_1L_1 + 1}V_{1}}_{\text{Low Pass Character}} + \underbrace{\frac{s^2C_1L_1}{s^2C_1L_1 + 1}V_{2}}_{\text{High Pass Character}}
\end{align*}

\subsubsection{Analysis}
\begin{enumerate}
  \item \textbf{Poles (\(D(s) = 0\)):} \\
    For both low pass \(H_1(s)\) and high pass \(H_2(s)\):
    \begin{align*}
      s^2C_1L_1 + 1 &= 0 \\
      s^2C_1L_1 &= -1 \\
      s^2 &= -\frac{1}{C_1L_1} \\
      s &= -2 \pi f_p \\
      (-2 \pi f_p)^2 &= -\frac{1}{C_1L_1} \hspace{1cm} \text{[let \(s = -2 \pi f_p\)]} \\
      4 \pi^2 f_p^2 &= \frac{1}{C_1L_1} \\
      f_p &= \frac{1}{2 \pi \sqrt{C_1L_1}}
    \end{align*}
    Complex conjugate poles on the imaginary axis indicate
    a marginally stable second-order system.
  \item \textbf{Zeroes (\(N(s) = 0\)):} \\
    For low pass \(H_1(s)\):
    \begin{align*}
        N_1(s) = 0
        1 \neq 0
    \end{align*}
    No zeroes for \(H_1(s)\). \\
    For high pass \(H_2(s)\):
    \begin{align*}
        N_2(s) = 0 \\
        s^2C_1L_1 = 0 \\
        s_z &= 0 \hspace{1cm} \text{[let \(s = s_z\)]} \\
        f_z &= \frac{|s_z|}{2\pi} = 0 Hz
    \end{align*}
    Zero at the origin for \(H_2(s)\).
  \item \textbf{DC Gain \(A_{DC}\):} \\
    For low pass \(H_1(s)\):
    \begin{align*}
      H_1(0) &= \frac{1}{0^2C_1L_1 + 1} = 1
      A_{DC_1} = 1 = 0 dB
    \end{align*}
    Unity DC gain for \(H_1(s)\). \\
    For low pass \(H_2(s)\):
    \begin{align*}
      H_2(0) &= \frac{0^2C_1L_1}{0^2C_1L_1 + 1} = 0
      A_{DC_2} = 0 = -\infty dB
    \end{align*}
    Zero DC gain for \(H_2(s)\).
  \item \textbf{Gain Bandwidth Product (GBW):}
    For low pass \(H_1(s)\):
    \begin{align*}
      GBW &= |A_{DC_1} \cdot f_p| \\
          &= 1 \cdot \frac{1}{2 \pi \sqrt{C_1L_1}} \\
          &= \frac{1}{2 \pi \sqrt{C_1L_1}}
    \end{align*}
    For high pass \(H_2(s)\):
    \begin{align*}
      GBW &= |A_{DC_2} \cdot f_p| \\
          &= 0 \cdot \frac{1}{2 \pi \sqrt{C_1L_1}} \\
          &= 0
    \end{align*}
    The high pass filter has no gain at DC, so the GBW which
    is a quantity measured from DC to the cutoff frequency is
    zero.
  \item \textbf{Impedance:}
    Both low pass and high pass filters have the same input
    and output impedance.
    Substituting \(G_1 = \frac{1}{R_1}\) and \(s = j \omega\):
    \paragraph{At Input:}
      \begin{align*}
        Z_{in} &= sL_1 + \frac{1}{sC_1} \\
        Z_{in} &= j \omega L_1 - j \frac{1}{\omega C_1} \\
        Z_{in} = j \left( \omega L_1 - \frac{1}{\omega C_1} \right)
      \end{align*}
    \paragraph{At Output:}
      \begin{align*}
        Z_{out} &= \frac{1}{sC_1} || sL_1 \\
        Z_{out} &= \frac{\left( \frac{1}{sC_1} \cdot sL_1 \right)}{\left( \frac{1}{sC_1} + sL_1 \right)} \\
        Z_{out} &= \frac{\frac{L_1}{C_1}}{s^2 L_1 C_1 + 1} \\
        Z_{out} &= \frac{\frac{L_1}{C_1}}{(j \omega)^2 L_1 C_1 + 1} \\
        Z_{out} &= \frac{\frac{L_1}{C_1}}{- \omega^2 L_1 C_1 + 1} \\
        Z_{out} = \frac{\frac{L_1}{C_1}}{1 - \omega^2 L_1 C_1}
      \end{align*}
\end{enumerate}

\subsubsection{Question}
Design a LC Low Pass Filter with a cutoff frequency of
\(f_p = 1 kHz\) using standard component values.
\paragraph{Solution:}
\begin{enumerate}
  \item Choosing \(C_1 = 1 \mu F\) (standard value).
  \item Calculating \(L_1\):
    \begin{align*}
      f_p &= \frac{1}{2 \pi \sqrt{C_1L_1}} \\
      1 kHz &= \frac{1}{2 \pi \sqrt{C_1L_1}} \\
      \sqrt{C_1L_1} &= \frac{1}{2 \pi \cdot 1 kHz} \\
      C_1L_1 &= \left( \frac{1}{2 \pi \cdot 1 kHz} \right)^2 \\
      C_1L_1 &\approx 2.533 \times 10^{-8} \\
      L_1 &= \frac{2.533 \times 10^{-8}}{C_1} = \frac{2.533 \times 10^{-8}}{1 \mu F} = 25.330 mH
    \end{align*}
    Standard value: \(L_1 = 22 mH\)
  \item Verifying \(f_p\):
    \begin{align*}
      f_p &= \frac{1}{2 \pi \sqrt{C_1L_1}} \\
      f_p &= \frac{1}{2 \pi \sqrt{1 \mu F \cdot 22 mH}} \\
      f_p &\approx 1073 Hz \\
    \end{align*}
    Close to desired \(1 kHz\).
\end{enumerate}
Design a LC High Pass Filter with a cutoff frequency of
\(f_c = 10 kHz\). Choose standard component values for \(L_1\)
and \(C_1\).
\paragraph{Solution:}
\begin{enumerate}
  \item Choosing standard value \(C_1 = 10 nF\).
  \item Calculating $L_1$:
    \begin{align*}
      f_c &= \frac{1}{2 \pi \sqrt{C_1L_1}} \\
      L_1 &= \frac{1}{(2 \pi f_c)^2 C_1} \\
      L_1 &= \frac{1}{(2 \pi \cdot 10 \times 10^3)^2 \cdot 10 \times 10^{-9}} \\
      L_1 &\approx 25.33 mH
    \end{align*}
    Choose standard value \(L_1 = 27 mH\).
  \item Verifying $f_p$ with chosen values:
    \begin{align*}
      f_p &= \frac{1}{2 \pi \sqrt{C_1L_1}} \\
      f_p &= \frac{1}{2 \pi \sqrt{0.01 \times 10^{-6} \cdot 27 \times 10^{-3}}} \\
      f_p &\approx 9.685 kHz
    \end{align*}
    Close to desired pole frequency.
\end{enumerate}

